\scanpage{100}
dégénérée. On peut passer même à l'hensélisé (non strict) de $\underline{\mathcal O}_s$, on \unsure{trouve|dans} la situation d'un trait hensélien $S$, à corps résiduel $k$ (ici $\mathbf Q$ --- ou une extension finie de $\mathbf Q$), d'une courbe $U_\eta$ de type $(g,\nu)$ (ici $(1,1)$) sur son pt générique $\eta$, on s'intéresse à l'extension $\mathcal E_{U_\eta}$ de $\mathcal E_\eta$ par $\pi$, où ici $\mathcal E_\eta$ s'inscrit encore dans
\[
(30)\qquad 1\longrightarrow T\longrightarrow\mathcal E_\eta\longrightarrow\mathcal E_k\longrightarrow1
\]
\deleted{On est ici en \unsure{car.} $0$} \inserted{en plus de}
\[
(31)\qquad 1-\pi-\mathcal E_{U_\eta}\longrightarrow\mathcal E_\eta\longrightarrow1
\]
Ici l'action extérieure de $T$ sur $\pi$ n'est sûrement pas triviale (je conjecture qu'elle l'est ssi il y a ``bonne réduction'') --- donc le groupe $\mathcal E_k$ n'opère pas lui $=$ extérieurement sur $\pi$. Mais tout scindage de (30)
\margin{mais je l'ai démontré, dans le contexte des VA, il y a belle lurette ! \unsure{Ça doit bien} être connu sur $H_1$ sûrement \unsure{!!}}
\margin{99}
